\section{Summation, stopping, and convergence}\label{sec:closure}

We now remove the size restrictions used in the finite identities.
There are two distinct inductions. The numerical cutoffs are chosen
backwards, starting with the last layer. The estimates for events of
the actual process are proved forwards, starting with the first
layer. An event in layer $j$ is bounded by a finite witness whose
input is at the beginning of that layer; only layers strictly before
$j$ are expanded to estimate that input. In particular, no
propagation-of-chaos assertion is used to bound a stopping event.

All constants denoted by $C$ in the few-crossing calculation below
are independent of the number $L$ of layers. They may depend on
$T$, on the potential, on the fixed exponent $\beta _0$ of
Proposition~\ref{pre:envelopes}, and on the corresponding bounds for
$f$ and the initial data. Constants denoted by $C_L$ may also depend
on $L$. The compact scattering cutoffs used in
Propositions~\ref{hist:irregular} and \ref{exp:defects} will be
chosen after $L$; they do not enter the few-crossing constant $C$.

\subsection{Completed finite histories}

A \emph{completed history} has no unresolved stopped factor, has
been expanded through all the layers that it enters, and has
undergone the simple-family cancellations at its centered cuts.
Its inputs are independent copies of $f$, a one-body residual from
Proposition~\ref{cent:residual}, or an initial centered factor from
Proposition~\ref{pre:initial}.

We shall need summability also when the highest expanded layer is left
uncentered. There are three applications: the finite reference
cutoff mismatch $B_j$ in \eqref{cent:reference-mismatch}; a
positive history ending in the terminal witness of
Definition~\ref{def:cell-terminal-witness}; and the completed
cells above a physical stop, together with its positive
whole-component continuation. In each case we fix the finite
upper block and expand only its lower centered inputs through
earlier layers.

An \emph{admissible mixed uncentered upper block} lies within one,
possibly shortened, layer. It consists of the signed
completed whole-cell terms of Section~\ref{sec:cell}, any
independent reference families and residual upper slots from
Section~\ref{sec:centering}, and, when a physical stop or a witness
requires it, a positive actual-component prefix and an optional
final trial cell. All original force systems, contact indicators,
pattern restrictions, and whole-component cut rules are retained.
Only its lower centered inputs are expanded through earlier
layers, with the usual simple-family cancellations there. No
cancellation of the upper block's roots is required.

The starting top labels are counted once. A reference root reuses
its deleted upper label and a residual is a one-body source on
that label. Choosing physical versus deleted-reference top slots
costs at most $2^{h_0}$. Bounded residual and sign alternatives
cost at most another $C^{h_0}$, included in the baseline $C^p$.
At a birth its original parent fixes its physical or reference
family. There is no additional per-birth family choice, and no
independent choice of a layer for an already chronologically
numbered birth.

At the final endpoint of a one-root witness, ``root'' means a
distinguished label and spatial anchor only. All particles of the
final trial system terminate together there, with their full
component Hamiltonians retained. No singleton-isolation condition
is imposed on the distinguished label at that endpoint. The
positive-observable variant in
Definition~\ref{def:cell-terminal-witness} and
the terminal convention in Lemma~\ref{hist:gaussian} include
this case. Earlier cuts obey the usual whole-component rules.

Write
\[
 p=h_0+n,\qquad n=\sum_jn_j,\qquad
 q_j=h_j+n_j,\qquad H=\sum_jh_j.
\]
Here $p$ counts independent formal label lifetimes once, $h_0$ is
the number of starting upper roots, $n_j$ is the number of labels
added in layer $j$, and $h_j$ is the number of roots entering that
layer. Starting roots may be centered; after a restart they need
not be. In the charge inequality we use $h_0$ as an upper bound
for the number of uncentered restarting roots. Thus $H$ counts
entries into expanded layers, including the first one, but not
the initial inputs at time zero. Those inputs are the initial
charges in Proposition~\ref{cent:charges}. Every root position
in this section is integrated. A bounded test of those positions
and velocities is permitted.

\begin{lemma}[Retained-slot entropy]\label{close:entropy}
Suppose a term retains $k_j$ slots among $q_j$ available slots at
successive cuts. Put $Q=\sum_jq_j$ and $H'=\sum_jk_j$. Then
\begin{equation}\label{close:slot-bound}
 \prod_j\binom{q_j}{k_j}
 \leq\binom{Q}{H'}
 \leq\exp\left\{H'\log\frac{eQ}{H'}\right\}
 \quad (H'>0).
\end{equation}
In the history expansion $Q\leq p+C H$ and $H'\leq C H$, up to a
fixed number of choices per retained slot. Consequently, when $H\leq\delta p$,
the logarithmic selection cost is at most
\begin{equation}\label{close:entropy-small}
 C\delta p\left(1+\log\frac1\delta\right).
\end{equation}
For general $H\leq Lp$, its cost is at most $C_L^p$.
\end{lemma}
\begin{proof}
The coefficient of $z^{H'}$ in
$\prod_j(1+z)^{q_j}=(1+z)^Q$ is the sum of the products on the
left of \eqref{close:slot-bound} over all allocations of $H'$.
The first inequality follows. The second follows from
$\binom Qk\leq(eQ/k)^k$. A label contributes an available slot at
its first input cut, and an additional available slot each time it
is retained across another cut. This gives the asserted bound on
$Q$. A bounded number of choices per retained slot contributes
$\exp(C H)$. Bounded choices at first-input slots instead cost
$C^p$; these belong to the baseline constant in the history
estimate, not to its small retained-slot entropy.
The last two assertions follow from the displayed inequalities,
using $\binom Qk\leq2^Q$ in the general case.
\end{proof}

\begin{lemma}[Energy and the short-layer factors]\label{close:energy}
For a completed history with the above counts, the parent, time,
and velocity factors in Proposition~\ref{hist:basic}, before
retained-slot and spatial losses, have total integral at most
\begin{equation}\label{close:energy-factor}
 (C b\sqrt L)^p
 \exp\{H+h_0\log_+(1/b)\},
 \qquad b\leq1.
\end{equation}
The constant $C$ is independent of $L$. A fixed part of the source
Gaussian remains available for the spatial estimates.
\end{lemma}
\begin{proof}
Use the source sum and augmented squared-speed budget of
Lemma~\ref{hist:gaussian}:
\[
 S=\sum_{\text{independent inputs }a}|v_a|^2,
 \qquad \mathcal E=S+2Bp.
\]
By Lemma~\ref{hist:gaussian}, every active or mixed sum of
squared speeds is at most $\mathcal E$. The same assertion bounds
the sum of squares of the velocities of the independently added
free rays. The source sum $S$ is not replaced by terminal
kinetic energy: a departing interacting component may have positive
potential energy. Retaining its full Hamiltonian gives instead
the Gaussian comparison
\begin{equation}\label{close:source-gaussian}
 e^{-\beta _0S}
 \leq e^{\beta _0Bp}
 e^{-(\beta _0/2)\sum_a|V_a|^2}
 e^{-(\beta _0/2)S},
\end{equation}
where $V$ is the vector of mixed terminal velocities. Further
splitting into a fixed number of portions is allowed. All these
exponents are fixed independently of the layer subdivision.

The source energy depends on the parent prescription. We therefore
first partition its domain into shells
\[
 rp\leq S<(r+1)p,\qquad r=0,1,\ldots.
\]
On shell $r$ use the deterministic upper budget
$E_r=(r+1+2B)p$. Parent choices that cannot belong to this
shell are killed. The parent-sum bound of
Proposition~\ref{hist:basic} then applies with this common budget;
it does not require moving a prescription-dependent energy through
a sum over parents. At a birth in layer $j$, if $u\leq q_j$ paths
are present, its summed flux is bounded by
\[
 C u\left(1+|w|+\sqrt{E_r/u}\right)
 \leq Cq_j\left(1+|w|+\sqrt{E_r/q_j}\right).
\]
Let $m$ count multiple-insertion participants and endpoint pins,
and let $\mathcal O$ be the ordinary births. The nonordinary
shell-speed factor in Proposition~\ref{hist:basic} is
$(r+2)^{m/2}$. Pad the speed product with $h_0$ unit factors.
AM--GM and $\sum_jn_j/q_j\leq L$ give
\begin{align}
 &(r+2)^{m/2}
 \prod_{i\in\mathcal O}
 \left(1+|w_i|^2+\frac{E_r}{q_{j(i)}}\right)^{1/2}
 \nonumber\\
 &\hspace{12mm}\leq
 \left(C L(r+2)+\frac{|V|^2}{p}\right)^{p/2}.
 \label{close:amgm}
\end{align}
Indeed the sum of the $p$ squared factors is at most
$p+m(r+1)+|V|^2+LE_r$.
This form of the estimate remains valid after the parent
majorization has enlarged the velocity integral to the full
product Gaussian; it does not incorrectly retain a source-shell
constraint on that enlarged integration domain.

For a fixed product Gaussian $\gamma$,
\[
 \int\left(1+\frac{|V|^2}{p}\right)^{p/2}\,d\gamma(V)
 \leq C^p.
\]
To check this, reserve half of its Gaussian exponent and bound
$e^{-c u}(1+u/p)^{p/2}$ by $C^p$ for $u\geq0$; the remaining
Gaussian integrates with an exponential-in-$p$ normalization.
Applying this to \eqref{close:amgm} gives
$[C\sqrt{L(r+2)}]^p$. A fixed portion of the source-energy
factor in \eqref{close:source-gaussian} gives $e^{-crp}$.
Thus the sum of the shell bounds is at most
\begin{equation}\label{close:shell-sum}
 C^p\sum_{r\geq0}e^{-crp}
 \bigl(C L(r+2)\bigr)^{p/2}
 \leq (C\sqrt L)^p.
\end{equation}
For the last inequality, bound
$e^{-cr/2}\sqrt{1+C(r+1)}$ uniformly in $r$ and sum the remaining
geometric series. Reserve an additional fixed Gaussian portion
before this calculation if spatial factors are present.

In each layer the ordered-time volume is $b^{n_j}/n_j!$. Moreover,
\[
 \frac{q_j^{n_j}}{n_j!}\leq e^{n_j+h_j},
\]
because $n!\geq(n/e)^n$ and
$n\log(1+h/n)\leq h$. Hence the parent and time product is at most
$b^n e^{n+H}$. Combining this with \eqref{close:shell-sum}, and
using $n=p-h_0$, proves \eqref{close:energy-factor}.
\end{proof}

There are two further summation points. First, the number of
nonnegative layer occupancies with total $n$ is
\begin{equation}\label{close:layer-count}
 \binom{n+L-1}{L-1}=\exp(o(p))
\end{equation}
for fixed $L$ and growing $p$. Births are chronologically ordered
within consecutive layer blocks; there is no independent choice
among $L$ layers for each already numbered birth. Second, the
spatial factors in Proposition~\ref{hist:basic}, allowing even a
separate localization at each cut, have logarithm at most
\begin{equation}\label{close:spatial-loss}
 C H+C H\log\left(1+\frac{L\mathcal E}{H}\right),
\end{equation}
with $H$ replaced by $H+h_0$ if the upper-cut convention requires
it. This follows by the weighted concavity of the logarithm
applied to the displayed top-line displacement bound in that
proposition. On shell $r$, a reserved factor $e^{-crp}$ absorbs
the dependence on $r$ in \eqref{close:spatial-loss}. If
$H\leq\delta p$, the remaining logarithmic cost is
\begin{equation}\label{close:spatial-small}
 C\delta p\left(1+\log\left(1+\frac L\delta\right)\right).
\end{equation}
For example, use
$1+a(r+c)\leq(1+ac)(1+r/c)$ with $c>0$ fixed, and reserve part
of $e^{-crp}$ for $\delta p\log(1+r/c)$. This also explains why
no independent Gaussian is needed for a carried root at each cut.

The multiple-insertion and endpoint-overlap mesh factors in
Lemma~\ref{hist:pinning} and Proposition~\ref{hist:basic} may be bounded
by their ordinary birth factors after $J$ is chosen. Indeed their
extra factors are negative powers of $J$ or positive powers of
$\varepsilon$, whereas all their remaining per-participant costs
are fixed powers of the polynomial label cap. Choosing the power
in $J$ last makes each such extra product at most one. This step
absorbs their constants, including constants depending on $L$;
it does not introduce a new $C_L$ for every ordinary single birth.
Their unbounded shell-energy factors have already been included
in \eqref{close:amgm}, rather than absorbed by a fixed mesh power.

\begin{proposition}[Summability of completed finite terms]
\label{close:summability}
Fix any desired number $a>0$. One can choose a finite $L$, then
the rarity cutoffs, with the following property. For every fixed
polynomial label cap $p\leq P^A$, choose $J$ to be a sufficiently
large fixed power of $P$. There exist $\eta_L>0$ and $p_L<\infty$
such that, for sufficiently small $\varepsilon$, the sum of the
integrated absolute weights of completed histories with exactly
$p$ labels is at most
\begin{equation}\label{close:complete-tail}
 e^{-ap},\qquad p\geq p_L,\qquad h_0\leq\eta_Lp.
\end{equation}
The same bound holds for an admissible mixed uncentered upper
block followed by centered lower layers, after the lower
simple-family cancellations. It is uniform over shortened endpoint
cells and layers and over tests of absolute value at most one.
For the original centered expansion, in which the upper
simple-family cancellations are also made, every fixed-size term
with a retained centered upper root tends to zero after integration.
\end{proposition}
\begin{proof}
First choose $L$ so that $b=T/L$ is within the short interval of
Proposition~\ref{cent:residual}, is at most one, and satisfies
\[
 C T/\sqrt L<e^{-4a-20}.
\]
Here $C$ is the constant in Lemma~\ref{close:energy}, including
the independent source-envelope constants. Its independence of
$L$ is essential. Next choose $\delta>0$ small enough that
\eqref{close:entropy-small}, \eqref{close:spatial-small}, and the
factor $e^H$ together cost at most $e^{ap/4}$ whenever
$H\leq\delta p$. Choose $\eta_L$ still smaller so that the
upper-root factor in \eqref{close:energy-factor} and any
upper-root localization also have logarithm at most $ap/4$ for
$h_0\leq\eta_Lp$. The finite occupancy factors in
\eqref{close:layer-count} are absorbed by increasing $p_L$.
The histories with few crossings, $H\leq\delta p$, then obey
\eqref{close:complete-tail} with room to spare.

For the complementary histories apply
Proposition~\ref{cent:charges}, after reducing $\eta_L$ to at
most $\delta/(2L)$. If $d$ is the number of distinct marked
lives, $r$ the number of residual factors, and $i$ the number of
retained initial singleton inputs, then
\[
 d+r+i\geq\theta_Lp,\qquad
 \theta_L=\frac{\delta}{2L^2}>0.
\]
The fixed-factor transfer between algebraic families and the
geometric pin forest in Remark~\ref{cent:pin-forest} can be
absorbed by reducing $\theta_L$. A new geometric first pin is
still a single birth. If it reroutes a birth to another displayed
family, the original required later encounter is a defect; it is
not counted as another birth.

The unrestricted finite-term bound is $C_L^p$ by
Proposition~\ref{hist:basic}. If at least $\theta_Lp/3$ charges
are residual or initial factors, their additional weight is at
most
\[
 \gamma_\varepsilon^{\theta_Lp/3},\qquad
 \gamma_\varepsilon=
 \max\{C\varepsilon^c,C\eta_\varepsilon\}
 \longrightarrow0,
\]
where $\eta_\varepsilon$ is the supremum of the fixed summed
spatial Gaussian residual norms in
Proposition~\ref{cent:residual}, taking the maximum over the
finitely many choices $K_j$. Its common exponent is not weakened
through the layers. This contribution beats any fixed
$C_L^p$ once $\varepsilon$ is small. No rate for the residual is
needed.

Otherwise an extensive number of lives is marked. The tunable
parameter, tree, and multiple-insertion discards of
Proposition~\ref{hist:irregular} and
Lemma~\ref{exp:terminal}, followed by
Proposition~\ref{exp:defects}, give a bound of the form
\begin{equation}\label{close:many-cross}
 C_L^p e^{-A_1p}
 +P^{C_Lp}\varepsilon^{c_Lp}.
\end{equation}
The desired $A_1$ is chosen after $L$ and $\theta_L$. Choose it
larger than all the fixed exponential costs plus $4a+20$.
Reduce $\eta_L$ to the top-root fraction supplied by
Proposition~\ref{exp:defects} for affected fraction $\theta_L/3$
and this accuracy, and increase $p_L$ to its fixed size threshold.
These choices leave the preceding few-crossing estimates valid.
The second term is also at most $e^{-(4a+20)p}$ for small
$\varepsilon$, because $p\leq P^A$ and
$\log P=o(|\log\varepsilon|)$. The expensive exposure guesses occur
only in this geometric-gain term, not in the term controlled
solely by residual smallness. This proves
\eqref{close:complete-tail} in the many-crossing case.

The argument used only the finite-history estimates and the
finite centered identities. An admissible mixed upper block is
allowed because the charge estimate imposes no condition on its
$h_0$ upper roots. Signed completed cells, reference families,
and a positive actual prefix are all admissible finite terms;
taking their absolute weights does not change their force rules
or contact constraints. The upper physical/reference choices
cost only $2^{h_0}$, with bounded alternatives absorbed in $C^p$.
The original parent determines a birth's factor, so these choices
introduce no extra parent power. Ordinary births across all
simultaneous factors retain the ordered-time bound; multiple
insertions have the mesh factors already estimated. Below this
one upper layer every retained input is an ordinary centered
slot, so Proposition~\ref{cent:charges} applies with precisely
the same allowance $Lh_0$. The source, spatial, and parent
bounds, and hence both crossing cases, remain unchanged.

Finally fix the total number of labels. The path argument in
Proposition~\ref{cent:charges} makes every surviving centered
upper root lead to an initial small factor, a vanishing residual,
or a marked finite history. For the last case, first restrict
velocities and scattering parameters to compact regular sets.
Corollary~\ref{exp:finite} makes its integral tend to zero. The complement has arbitrarily small integral by the
Gaussian and flux bounds, so the compact restriction may then
be removed. This finite-size argument does not mistake a fixed
small-parameter cutoff for a quantity tending to zero. Ordered
time domination and shortened cells make the argument uniform
in the endpoint time.
\end{proof}

\begin{remark}[Dependence on the polynomial size cap]
The rarity cutoffs, top-root fraction, and fixed size threshold in
Proposition~\ref{close:summability} can be chosen independently of
the exponent $A$ in $p\leq P^A$.
Indeed, the source-energy cutoff follows from
$C_L^p e^{-cRp}$, the compact-parameter cutoffs from their one-step
probabilities, and the terminal-tree cutoff from
$C_*^p M^{-\zeta\alpha p/16}$, with affected fraction $\alpha$ and
fixed comparison exponent $\zeta>0$ as in Lemma~\ref{exp:terminal}.
These choices depend on the fixed layer count, source bounds, affected
fraction, and requested exponential accuracy, but not on $A$.
In the geometric estimates, the losses involving $p$ and $J=\lceil P^Q\rceil$ are
polynomial powers paid by positive powers of $\varepsilon$.
For each fixed $A,Q$ and $c>0$,
\[
 P^{C(A,Q)p}\varepsilon^{cp}
 =\exp\!\left\{p\bigl(C(A,Q)\log P-c|\log\varepsilon|\bigr)\right\}
\]
is smaller than any prescribed exponential in $p$ for sufficiently
small $\varepsilon$. For a fixed crowding fraction $\delta$, the grid
loss is absorbed by the fixed threshold $p\geq8/\delta$.
For the multiple-insertion mesh factor,
choose $Q>2dA$ as in Proposition~\ref{hist:basic}.
Thus only the required mesh exponent and the sufficiently-small-$\varepsilon$
threshold depend on the polynomial cap. This permits that cap to be
calculated from the deterministic stopping schedule after the rarity
cutoffs have been fixed, and the cell mesh to be chosen last.
\end{remark}

\subsection{Finite witnesses for actual-process events}
\label{close:witnesses}

We give an auxiliary event slightly stronger than separate caps
on trial components and actual backward histories. It makes the
inherited-threshold issue explicit.

Fix a full mesh of $[0,T]$ with layers $1,\ldots,L$ and cells
of length $\Delta=b/J$. At the start of a layer, the ancestry of
a label is its entire instantaneous polymer. At the end of a
cell, define its \emph{actual layer ancestry} recursively as the
union of the ancestries, at the start of that cell, of all labels
in its actual whole-cell contact component. This definition uses
the full Newtonian dynamics and resets at each layer start.
Labels in one instantaneous polymer at a mesh cut have the same
ancestry: they belonged to the same preceding whole-cell
component, or to the same initial polymer at the layer start.

At a cell start take any trial subsystem that is a union of the
actual instantaneous polymers. Evolve that subsystem in
isolation during the cell. For each component of its cumulative
contact graph, take the union of the actual layer ancestries of
its labels at the cell start. Its \emph{ancestral size} is the
cardinality of that union. For a threshold $D$, let
$\mathcal B_j(D)$ be the event that an ancestral size exceeds
$D$ somewhere in the first $j$ layers, or that a layer-start
polymer there already exceeds $D$. All trial subsystems are
included in this definition. In particular, on its complement
every trial cell component has at most $D$ labels, and every
actual backward history within a layer has at most $D$ labels.

\begin{lemma}[First witness at an inherited threshold]
\label{close:witness-extraction}
Let $D\geq1$ be an integer. For a finite microscopic cloud, a first occurrence of
$\mathcal B_j(D)$ admits a witness in some layer $i\leq j$
with one distinguished integrated upper label and between $D$
and $2D^2$ formal labels. The witness consists of positive actual
component histories before its last cell and one trial subsystem
in its last cell. Every starting or departing set respects whole
instantaneous polymers. The bound is independent of the cutoff
assigned separately to any earlier layer. The weighted ancestral
construction above in fact permits $D<p\leq2D$.
\end{lemma}
\begin{proof}
Throughout this proof the threshold is the same number $D$.
Take the first violating cell in the preceding prefix. There
cannot be a previously unaccounted giant polymer at its start.
A giant polymer at a layer boundary was contained in the
preceding actual cell component, and hence would have violated
the same threshold earlier. Initially the exclusion makes all
polymers singletons.

Before the violating cell, every current polymer has an actual
ancestry of size at most $D$. Regard these polymers as atoms,
with their ancestry sets as weights. In a violating trial system
follow the increasing contact-component graph until a union of
weights first exceeds $D$. Before this merger all participating
components have weight at most $D$. At a binary merger their
union has size at most $2D$. If several contacts occur together,
select connected pre-contact components one at a time, stopping
when the union first exceeds $D$. The last added component has
weight at most $D$, so the same bound holds. Thus ties require
no probabilistic convention. Before the selected merger these
components have had no contacts with the other trial components;
their trajectories up to that time are consequently their own
isolated trajectories.

Let $U$ be the union of the selected current polymers and let
$W$ be the union of their actual ancestries. Then
\[
 D<|W|\leq2D,\qquad U\subseteq W.
\]
Trace the members of $U$ backwards through their actual
whole-cell components to the beginning of layer $i$. The union
of these histories has exactly the required ancestry labels in
$W$. An ancestor not in $U$ terminates at its prescribed earlier
cut. In particular, it is not kept artificially active in the
last trial cell. The last cell runs only the isolated flow on
$U$. At every such removal a whole close component departs;
otherwise a contact across that cut would have included the
departing label in the next required polymer.

Extend this last isolated trial flow to the fixed end of its
cell. It already has the required connected-contact event before
the selected merger, even if its subsequent path differs from
the original larger trial system. All coordinate changes can
therefore be made at fixed times; no stopping-time Jacobian is
introduced. Choose a label $a\in U$ as the upper distinguished
label. The trial contact graph connects $U$ to this label, and
the actual ancestry graphs connect all remaining labels in $W$
to $U$ backwards, cell by cell. All of $U$ terminates at this
final endpoint. Its distinguished label need not be isolated
from $U\setminus\{a\}$ there; in particular, this construction
does not exclude a cluster persisting to the end of the cell.
Keep the full terminal Hamiltonian before applying stability,
and integrate all terminal coordinates. Connectivity supplies
$p-1$ translation pins, with one spatial anchor and hence
$h_0=1$; endpoint close components use the endpoint pins of
Lemma~\ref{hist:pinning}. This gives the asserted positive
history. Keeping the cruder union of at most $2D$ histories of
size at most $D$ would give the stated safe bound $2D^2$.
\end{proof}

The lemma is used as a union bound, not as a statement about a
conditioned law. The prior-good and first-occurrence restrictions
are used to extract the bounded witness and are then dropped.
Its start set is a union of actual polymers, so factorial
expectation integrates its nonnegative kernel against the
\emph{unconditional} isolated marginal at the start of layer
$i$. Only its last cell is a trial cell. There is no factor for
choosing an arbitrary subset of all $N$ polymers: the extracted
witness already records the selected labels, and the factorial
sum counts those labels once.

Index the witness sum by the resulting positive whole-cell
prescriptions, counting each prescription once. For a fixed
final set $U$ and distinguished label $a$, its last-cell kernel
is $c_I^{\{a\}}(U)\mathsf J^I_{\{a\},U}$, with any retained
witness restriction included as an indicator bounded by one.
The earlier cells are the successive positive root-component
decompositions described after Definition~\ref{cell:histories}.
Different extraction choices that yield the same prescription
do not produce additional summands.

Writing $\mathcal W_{i,m,D}$ for the sum of these positive
one-root witness kernels in cell $m$, the preceding observation
gives the finite inequality
\begin{equation}\label{close:witness-union}
 \mathbb P\bigl(\mathcal B_j(D)\bigr)
 \leq \mu\sum_{i\leq j}\sum_{m\leq J}
 \left\langle\mathcal W_{i,m,D},
       \widetilde F^\varepsilon(t_{i-1})\right\rangle .
\end{equation}
The notation on the right sums the relevant input orders between
$D+1$ and $2D^2$ with their factorial weights. The factor $\mu$
undoes the normalization of the single distinguished label.
Birth-cell choices inside each witness are summed by ordered
contact times as in Proposition~\ref{hist:basic}; they introduce
no additional empty-cell multiplicity. The displayed $LJ$
choices locate only the final witnessing cell.

\subsection{Raw amplification and the cutoff schedule}

\begin{lemma}[Integrated amplification and bad clouds]
\label{close:amplification}
Consider a finite completed block containing at most $M$ labels
and acting on an input signed array in integrated total variation.
For $M$ and $J$ bounded by fixed powers of $P$, its operator norm
is at most
\begin{equation}\label{close:raw-amplification}
 \exp(C_LMP).
\end{equation}
The same bound holds for the pattern-summed centered norm of
\eqref{cent:centering-norm}. If a physical layer is stopped below
$K$ labels, the contribution of a
microscopic event $\mathcal A$ to the exact observable minus its
completed physical below-$K$ expansion is at most
\begin{equation}\label{close:bad-discard}
 \exp(C_LKP)\,\mathbb P(\mathcal A)^{1/2},
\end{equation}
provided its upper root count is less than $K$. This discard is
made before estimating any overshoot.
\end{lemma}
\begin{proof}
For a fixed term, product Hamiltonian transport preserves phase
volume, marginalization decreases total variation, and every
contact or isolation indicator has absolute value at most one.
Tensoring with an independent input has bounded integrated norm.
The factorial normalization costs at most $\mu^M$. With at most
$M$ labels, the finite subset, generation, cut, and parent
descriptions cost at most $(C_L M J)^{C_LM}$. This proves
\eqref{close:raw-amplification}, since
$\log M+\log J=O(\log P)=o(P)$.
The centering and inverse transforms have norm at most $2^M$
by \eqref{cent:centering-norm}; this factor is absorbed by
\eqref{close:raw-amplification}. Pattern summation introduces no
additional partition count, by Corollary~\ref{cent:pattern-collapse}.

For \eqref{close:bad-discard}, first work on an $N$-particle
cloud. The exact root observable is bounded by its root-tuple
count. Every completed below-$K$ term involves fewer than $K$
labels. Its absolute value, summed over their choices, is thus
bounded by $(C_LKJ)^{C_LK}(1+N)^K$, times normalization factors
no larger than $\mu^K$. No bound for the number of terms with
an unbounded last component is used here: we are bounding the
exact observable minus the completed small terms on $\mathcal A$.

The initial factorial bound gives
$\mathbb E(N)_\ell\leq\mu^\ell$, and $N$ is conserved. Expanding
ordinary moments in factorial moments therefore gives, for
integer $r\geq1$,
\[
 \mathbb E(1+N)^r\leq \bigl(Cr(1+\mu)\bigr)^r.
\]
Cauchy--Schwarz with $r=2K$, followed by
$\log\mu=2|\log\varepsilon|$, proves
\eqref{close:bad-discard}. These finite polynomial moments also
justify truncating first to $N\leq N_0$ and then taking
$N_0\to\infty$ in all the capped identities. The original
unconditioned ensemble is used throughout; no truncated ensemble
is recentered in place of it.
\end{proof}

Here is an explicit noncircular schedule. Choose $L$, the crossing
fraction, and the fixed rarity cutoffs as in
Proposition~\ref{close:summability}, with $a=32$. Choose a fixed
integer $\nu\geq4$ large enough to dominate the polynomial
description and witness-size bounds in the finite constructions.
Starting at layer $L$, put $U_L=\lceil P\rceil$. For
$j=L,L-1,\ldots,1$, define
\begin{align}
 K_j&=\lceil P^3U_j\rceil,&
 D_j&=\lceil P^3K_j\rceil, \label{close:cutoffs}\\
 W_j&=\left\lceil(1+U_j+K_j+D_j)^\nu\right\rceil. \nonumber
\end{align}
When $j>1$, choose
\begin{equation}\label{close:upper-cap}
 U_{j-1}=
 \left\lceil\left(1+U_j+W_j+
                 \sum_{\ell>j}W_\ell\right)^\nu\right\rceil.
\end{equation}
The symbols $U_j$ are caps for the entire already constructed
upper block, not just for its currently retained roots. Increasing
$\nu$ once, if necessary, covers every finite composition of the
stated deterministic polynomial bounds. There are only $L$
levels. All these quantities are fixed powers of $P$ up to
constant factors. In particular,
\begin{equation}\label{close:hierarchy}
 \frac{K_j}{U_jP}\longrightarrow\infty,\qquad
 \frac{D_j}{K_jP}\longrightarrow\infty,
 \qquad U_i\geq W_j\quad(i<j).
\end{equation}
Every possible inherited $D_j$-witness in an earlier layer is
bounded by a polynomial in this same $D_j$, and is included in
the downstream demands in \eqref{close:upper-cap}. We never
require $K_i$ to dominate its own $D_i$.

Choose a fixed $A_*$ with every complete or stopped capped
history containing at most $P^{A_*}$ labels. Finally take
$J=\lceil P^{A_{\rm cell}}\rceil$, with $A_{\rm cell}$ larger
than all the fixed polynomial exponents required by the estimates in
Propositions~\ref{hist:basic}, \ref{hist:irregular},
\ref{hist:crowding}, and~\ref{exp:defects}.
Increasing $A_{\rm cell}$ affects raw
amplifications only through $O(\log P)$, so it does not change
\eqref{close:hierarchy}. This is why the cell mesh is chosen last.

\subsection{Chronological removal of the stopped terms}

For each layer distinguish the two exact remainders in
\eqref{cent:layer-identity}. The physical stopping identity
\eqref{cent:physical-layer} has remainder
$\mathsf R_j^{\mathrm{phys}}$. After centering and the exact
replacement $f=\mathscr G+r$, the finite high-total-size part
$B_j$ of \eqref{cent:reference-mismatch} is also present. Thus
\begin{equation}\label{close:centered-remainder}
 E(t_j)=\mathsf A_jE(t_{j-1})+\mathsf Z_j,\qquad
 \mathsf Z_j=
 \mathsf M_{t_j}\mathsf R_j^{\mathrm{phys}}
       +B_j[E(t_{j-1})].
\end{equation}
The argument of $B_j$ is the true lower centered family.
This finite reference cutoff mismatch is not a stopped-cloud
functional and is not bounded by positive physical domination.

Use the physical and pattern-summed norms from
\eqref{cent:centering-norm}, taking a supremum also over shortened
ending times in layer $j$, and put
\begin{equation}\label{close:remainder-norms}
 \begin{split}
 Q_j&=\sup_t\|\mathsf R_j^{\mathrm{phys}}(t)\|_{1,\leq U_j},\\
 V_j&=\sup_t\|B_{j,t}[E(t_{j-1})]\|_{\pi,1,\leq U_j},\\
 Z_j&=\sup_t\|\mathsf Z_j(t)\|_{\pi,1,\leq U_j},\qquad
 R_j=\max\{Q_j,Z_j\}.
 \end{split}
\end{equation}
Here $B_{j,t}$ denotes the same finite operator with the layer
shortened at $t$. In particular,
\begin{equation}\label{close:centering-remainder-norm}
 Z_j\leq2^{U_j}Q_j+V_j.
\end{equation}
The scalar $R_j$ therefore controls both the physical interruption
and the complete centered remainder used in the lower-layer
recurrence. Physical remainders are aggregate signed differences,
not sums of the norms of unbounded stopped descriptions. On a
good cloud we use their positive domination; on a bad cloud we
use the exact difference before taking absolute values.

\begin{proposition}[Actual witnesses and layer remainders]
\label{close:tails}
With the preceding schedule, for sufficiently small
$\varepsilon$ and every $1\leq j\leq L$,
\begin{equation}\label{close:actual-tail-bounds}
 \mathbb P\bigl(\mathcal B_j(D_j)\bigr)\leq e^{-8D_j},
 \qquad R_j\leq e^{-8K_j}.
\end{equation}
The estimates use the original Newtonian process, without
conditioning it to stay in a good set.
\end{proposition}
\begin{proof}
The finite-cloud stopping and witness constructions are performed
first. Integrate their finite identities with $N\leq N_0$ and
then let $N_0\to\infty$, using
Lemma~\ref{close:amplification} and the finite expectations in
Proposition~\ref{cell:stopping}. Only after this passage do we
apply the centered identities to the original ensemble and its
initial insertion formula. In particular, no centering bound for
a law conditioned on $N\leq N_0$ is asserted.

We induct on $j$ in chronological order, with $R_k$ already
controlling both remainder types whenever $k<j$.
First estimate the finite reference cutoff mismatch. By
\eqref{cent:mismatch-size}, every term of $B_j$ has
\[
 K_j\leq p<(U_j+1)K_j,\qquad h_0\leq U_j.
\]
This cap is below $W_j$. Treat this finite product of signed
physical terms, reference families, and residual upper slots as
an uncentered mixed upper block. Expand its true input
$E(t_{j-1})$ by \eqref{cent:layer-identity} only through the
strictly earlier layers $j-1,\ldots,1$. Apply the usual
simple-family cancellations at those lower cuts. Every completed
history still contains at least $K_j$ labels, and its starting
top count is at most $U_j=o(K_j)$. Therefore
Proposition~\ref{close:summability} bounds their total norm by
$C_L e^{-32K_j}$. Every interruption in a layer $k<j$ contains
the complete centered remainder $\mathsf Z_k$ from
\eqref{close:centered-remainder}. Its upper block fits $U_k$
by \eqref{close:upper-cap}, so
Lemma~\ref{close:amplification} bounds it by
$e^{C_LU_kP}R_k$. We have proved
\begin{equation}\label{close:reference-cutoff-bound}
 V_j\leq C_L e^{-32K_j}
       +C_L\sum_{k<j}e^{C_LU_kP}R_k.
\end{equation}
No good-cloud event or physical overshoot is used in this
estimate. In particular, neither $R_j$ nor a bound on
$\mathcal B_j(D_j)$ occurs on its right side: $B_j$ is already
a finite upper block, and no layer-$j$ expansion is inserted
below it.

To prove the probability bound for $\mathcal B_j(D_j)$, apply
Lemma~\ref{close:witness-extraction} with the fixed threshold
$D_j$ throughout the preceding prefix. A witness whose last cell
is in layer $i\leq j$ has more than $D_j$ labels and at most
$2D_j^2$. Drop its first-occurrence restrictions, as in
\eqref{close:witness-union}. Expand its start marginal through
layers $1,\ldots,i-1$, using the exact finite centered recurrence
\eqref{cent:layer-identity}, including both remainder terms.
Every completed term has at least $D_j$
labels and one integrated uncentered upper root. Its total
integral is therefore at most $C e^{-32D_j}$ by
Proposition~\ref{close:summability}. This is a statement about
finite terms with independent inputs; it does not presume a
bound for the probability being proved.

Every interrupted lower term occurs in a layer $k<i$. Its
upper block is bounded by $U_k$ by construction. Its contribution
is at most $e^{C_LU_kP}R_k$ by
Lemma~\ref{close:amplification}, and $R_k$ is already bounded
by the induction hypothesis. Thus
\begin{equation}\label{close:prob-induction}
 \mathbb P\bigl(\mathcal B_j(D_j)\bigr)
 \leq \mu LJ\left(
 C_L e^{-32D_j}
 +C_L\sum_{k<j}e^{C_LU_kP}R_k\right).
\end{equation}
For $j=1$ the sum is empty and all input estimates are initial
ones. For $k<j$, \eqref{close:hierarchy} makes
$e^{C_LU_kP}R_k\leq e^{-7K_k}$ for small
$\varepsilon$. The same schedule makes each $K_k$ exceed every
fixed multiple of $D_j+P+\log J$. Since
$\log\mu=2|\log\varepsilon|$, the right side of
\eqref{close:prob-induction} is at most $e^{-8D_j}$.

Now estimate the physical remainder norm $Q_j$. On
$\mathcal B_j(D_j)$, discard the exact observable minus the
completed below-$K_j$ terms before estimating an overshoot.
Lemma~\ref{close:amplification} bounds this part by
\begin{equation}\label{close:own-bad-cost}
 e^{C_LK_jP}
 \mathbb P\bigl(\mathcal B_j(D_j)\bigr)^{1/2}.
\end{equation}
On the complementary cloud, all trial components have at most
$D_j$ labels and every earlier actual history within the layer
has at most $D_j$ labels. In the first stopped cell,
Proposition~\ref{cell:stopping} therefore bounds its selected
union by
\[
 \max\{U_jD_j,K_j-1+D_j\}.
\]
Its positive continuation through the earlier cells of the layer
adds at most a factor $D_j$ in label count. Together with the
already completed upper part of that layer, this is less than
$W_j$ in \eqref{close:cutoffs}. Every such history still has at
least $K_j$ labels: the positive continuation never removes a
label counted by the stopping rule. Its upper root count is at
most $U_j$, and $U_j/K_j\to0$.

The completed signed cells above the stopped cell and the
positive actual prefix below it form an admissible mixed upper
block. Drop the remaining good-set restrictions and expand its
input through the strictly earlier layers, using the full
centered recurrence there. Completed terms are bounded by
$C_L e^{-32K_j}$ using
Proposition~\ref{close:summability}; all their labels are below
the fixed polynomial cap. Earlier interruptions again cost at
most $e^{C_LU_kP}R_k$. Consequently
\begin{equation}\label{close:physical-remainder-induction}
 Q_j\leq C_L e^{-32K_j}
 +e^{C_LK_jP}
       \mathbb P\bigl(\mathcal B_j(D_j)\bigr)^{1/2}
 +C_L\sum_{k<j}e^{C_LU_kP}R_k.
\end{equation}
Polynomially many final-cell and stopping choices are absorbed
by the exponential margins in this inequality.

Finally combine \eqref{close:reference-cutoff-bound},
\eqref{close:physical-remainder-induction}, and
\eqref{close:centering-remainder-norm}. Since $U_j=o(K_j)$,
the factor $2^{U_j}$ changes $e^{-32K_j}$ to at most
$e^{-31K_j}$ for small $\varepsilon$. Its multiplication of
the other terms is absorbed by increasing $C_L$ in their raw
amplifications. Thus the common scalar obeys
\begin{equation}\label{close:remainder-induction}
 R_j\leq C_L e^{-31K_j}
 +e^{C_LK_jP}
       \mathbb P\bigl(\mathcal B_j(D_j)\bigr)^{1/2}
 +C_L\sum_{k<j}e^{C_LU_kP}R_k.
\end{equation}
By the probability bound just proved and
\eqref{close:hierarchy}, its middle term is at most
$e^{-3D_j}$. Every earlier term is at most
$e^{-7K_k}$, and the earlier $K_k$ dominate $K_j$. The right
side is at most $e^{-8K_j}$ for small $\varepsilon$.

This completes the induction. Notice its order: the reference
cutoff mismatch and the probability of a layer-$j$ witness used
only remainders with index less than $j$; the physical
layer-$j$ remainder then used that newly proved probability.
Centering its physical remainder and adding the separately
estimated finite mismatch produced $\mathsf Z_j$ and the common
bound $R_j$. No estimate assumed its own layer's remainder.
The larger threshold $D_i$ assigned to a preceding
layer was never used to infer that a polymer exceeding $D_j$
was absent. The same-$D_j$ first-witness construction is what
justifies that inference.
\end{proof}

\subsection{Fixed-order chaos and the empirical consequence}

\begin{proof}[Proof of Theorem~\ref{thm:main} and Corollary~\ref{cor:empirical}]
Fix $s$. For sufficiently small $\varepsilon$, it is at most
$U_L$. Expand a retained centered singleton array of this order
with the exact finite identity \eqref{cent:finite-layers}.
Proposition~\ref{close:tails} and
\eqref{close:raw-amplification} remove both the physical and
the reference cutoff remainders:
the total error is bounded by a finite sum of
$e^{C_LU_jP}R_j$, which tends to zero.

For the completed terms, first fix an integer $m$ larger than
the finite constants required in
Proposition~\ref{close:summability}. The sum of all terms with
more than $m$ labels is at most $C_{s,L}e^{-m}$, uniformly for
small $\varepsilon$. The finitely sized part, with at most $m$
labels, tends to zero by the last assertion of that proposition.
Here finiteness refers to the number of labels and tree shapes;
the contact times are integrated over their ordered simplices.
Letting first $\varepsilon\to0$ and then $m\to\infty$ proves
integrated convergence of every fixed retained centered
singleton array.

These estimates are uniform over deterministic observation
times. Use the fixed full mesh of $[0,T]$ and truncate only its
last cell at the observation time. Full-cell trial-component
and ancestry caps imply the corresponding caps on every prefix
of that cell. All ordered-time estimates use intervals of
length at most $b$ and remain valid after shortening. The
residual and finite-size convergence are uniform for these
endpoints by Proposition~\ref{cent:residual}. At $t=0$ the
expansion is empty and Proposition~\ref{pre:initial} supplies
the conclusion directly. Thus no positive lower bound on the
observation time is needed.

Let
\[
 \mathcal D_s^\varepsilon=
 \{Z_s:|x_i-x_j|>\varepsilon\text{ for every }i\ne j\},
 \qquad
 G_s^\varepsilon=
 \mathbf1_{\mathcal D_s^\varepsilon}
       \widetilde F_s^\varepsilon.
\]
The inverse centered identity in
Proposition~\ref{cent:identity}, with the pattern restrictions
restored, gives
\begin{equation}\label{close:isolated-limit}
 \sup_{0\leq t\leq T}
 \|G_s^\varepsilon(t)-f(t)^{\otimes s}\|_1\longrightarrow0.
\end{equation}
Indeed every nonempty retained subset contributes one of the
centered arrays just bounded. Integrating the other independent
$f$ slots costs at most their mass one; the separation indicators
only decrease this bound. The missing independent mass outside
$\mathcal D_s^\varepsilon$ is $O_s(\varepsilon^3)$, uniformly
in time, because the spatial density of $f$ is bounded by its
uniform Gaussian envelope.

We have $0\leq G_s^\varepsilon\leq F_s^\varepsilon$. Since
$N$ is conserved,
\[
 m_s^\varepsilon:=\int F_s^\varepsilon(t)
 =\mu^{-s}\mathbb E(N)_s
 =\int F_s^\varepsilon(0)\longrightarrow1
\]
by Proposition~\ref{pre:initial}. Hence
\begin{align*}
 \|F_s^\varepsilon-f^{\otimes s}\|_1
 &\leq\|G_s^\varepsilon-f^{\otimes s}\|_1
       +\int(F_s^\varepsilon-G_s^\varepsilon)\\
 &\leq2\|G_s^\varepsilon-f^{\otimes s}\|_1
       +|m_s^\varepsilon-1|.
\end{align*}
Together with \eqref{close:isolated-limit}, this proves
\eqref{intro:chaos}. Thus close configurations, including
any bound clusters formed by the actual dynamics, are removed
by their nonnegative mass. They have never been excluded from
the microscopic evolution.

Finally put
\[
 A_\varepsilon(t)=\mu^{-1}\sum_{i=1}^N a(z_i(t)),
 \qquad a_f(t)=\int a(z)f(t,z)\,dz.
\]
Factorial normalization gives
\[
 \mathbb E A_\varepsilon(t)=\int aF_1^\varepsilon(t),
 \qquad
 \mathbb E A_\varepsilon(t)^2
 =\int a(z_1)a(z_2)F_2^\varepsilon(t)
   +\mu^{-1}\int a(z)^2F_1^\varepsilon(t,z)\,dz.
\]
It follows that
\begin{align*}
 \mathbb E|A_\varepsilon(t)-a_f(t)|^2
 \leq\|a\|_\infty^2\bigl(&
 \|F_2^\varepsilon(t)-f(t)^{\otimes2}\|_1
 +2\|F_1^\varepsilon(t)-f(t)\|_1
 +\mu^{-1}m_1^\varepsilon\bigr).
\end{align*}
The right side tends to zero uniformly over deterministic $t$.
Chebyshev's inequality proves \eqref{intro:empirical}.
This conclusion is the stated uniformity of probabilities at
deterministic times; it makes no assertion about a supremum
inside a sample path.
\end{proof}
